\section{Discussion}
\label{sec:discussion}

\paragraph{What the memory advantage was.}
The LSTM's 25 to 34-point advantage over a small memory-free MLP survives every statistical check
we applied, and it is real in the sense that it reproduces across seeds and maps. It is not,
for the most part, an advantage of remembering what was seen. Width alone explains about a tenth
of it, the time step about four fifths, and what is left for the LSTM beyond a clock is about
3 points on new maps, a difference we can detect but that is far smaller than the one we set out
to explain. The LSTM does need its recurrent state (erasing it costs 28 points), and its state
does encode the time step, and handing the time step to the same network with its state erased
gives back 24 of those 28 points. On this task, the recurrent state is used largely as a clock. Whether the remaining 3 points are memory of the scene or a
better clock is an open question.

\paragraph{Recommendations.}
For memory comparisons in tasks with time limits or time-dependent rewards:
\begin{enumerate}
  \item Give memory-free baselines the time step, or remove the dependence on time, before
  attributing a gap to memory. In our task this one input moved a memory-free network from 65\%
  to 92\%.
  \item Match parameters, and include an ablation that keeps the recurrent agent's weights but
  erases its state; together they separate capacity from recurrence.
  \item Report final success over paired seeds, with intervals: about ten paired seeds detect a
  10-point gap here, and about forty a 5-point gap. Avoid binary acquisition verdicts, or report
  them with many seeds and several evaluations per checkpoint.
  \item Re-evaluate saved checkpoints with new sampling seeds. It is cheap, it measures the noise
  of the evaluation directly, and it showed that the acquisition verdict of more than a quarter of
our LSTM runs (8 of 30) flips under new sampling seeds.
  \item Keep a sealed panel of new maps for the final comparison, and use more maps rather than
  more seeds when the claim is about generalization to new maps.
\end{enumerate}

\paragraph{Variance in procedurally generated environments.}
The decomposition shows why a single number such as ``use ten seeds'' cannot be right in general.
Final success varied mainly between seeds, so more seeds help and more evaluations help little;
the acquisition verdict varied mainly between evaluations, so the opposite holds; and claims about
unseen maps are limited by the number of maps, whatever the number of seeds. A short pilot with
re-evaluation is enough to estimate these components and to choose a design before the main
experiment.

\paragraph{Scope.}
Everything here is one task, one algorithm and small networks, and we expect the size of the
clock effect to depend on the horizon, on how strongly the reward depends on time and on how
value targets are computed at the time limit. The method, however, transfers: the same controls
and the same decomposition can be run on any memory benchmark at little cost, and we release the
tooling to do so.
